\section{From Neural Networks to R1CS: A Worked Example}

\subsection{Overview}

In zero-knowledge machine learning (zkML), we wish to prove that a model has been evaluated correctly without revealing inputs or parameters.

The key idea is:
\[
\text{model} \;\longrightarrow\; \text{arithmetic computation} \;\longrightarrow\; \text{constraint system (R1CS)}.
\]

A \textbf{Rank-1 Constraint System (R1CS)} consists of constraints of the form
\[
\langle A_i, w \rangle \cdot \langle B_i, w \rangle = \langle C_i, w \rangle,
\]
where $w$ is the witness vector.

\subsection{A Small Neural Network}

Consider a one-hidden-layer network:
\[
z = W_1 x + b_1, \quad h = \phi(z), \quad y = W_2 h + b_2,
\]
with
\[
W_1 =
\begin{bmatrix}
2 & -1 \\
1 & 3
\end{bmatrix}, \quad
b_1 =
\begin{bmatrix}
1 \\
0
\end{bmatrix},
\]
\[
W_2 =
\begin{bmatrix}
4 & -2
\end{bmatrix}, \quad
b_2 = 3.
\]

We choose a zk-friendly activation:
\[
\phi(t) = t^2.
\]

Thus the computation becomes:
\[
z_1 = 2x_1 - x_2 + 1, \quad
z_2 = x_1 + 3x_2,
\]
\[
h_1 = z_1^2, \quad h_2 = z_2^2,
\]
\[
y = 4h_1 - 2h_2 + 3.
\]

\subsection{Witness Vector}

Introduce variables for all intermediate values:
\[
w = (1, x_1, x_2, z_1, z_2, h_1, h_2, y).
\]

The leading $1$ is a constant wire.

\subsection{Constraint Construction}

Each computation step is converted into an R1CS constraint.

\paragraph{Constraint 1:}
\[
z_1 = 2x_1 - x_2 + 1
\quad \Rightarrow \quad
(1 + 2x_1 - x_2)\cdot 1 = z_1.
\]

\paragraph{Constraint 2:}
\[
z_2 = x_1 + 3x_2
\quad \Rightarrow \quad
(x_1 + 3x_2)\cdot 1 = z_2.
\]

\paragraph{Constraint 3:}
\[
h_1 = z_1^2
\quad \Rightarrow \quad
z_1 \cdot z_1 = h_1.
\]

\paragraph{Constraint 4:}
\[
h_2 = z_2^2
\quad \Rightarrow \quad
z_2 \cdot z_2 = h_2.
\]

\paragraph{Constraint 5:}
\[
y = 4h_1 - 2h_2 + 3
\quad \Rightarrow \quad
(3 + 4h_1 - 2h_2)\cdot 1 = y.
\]

\subsection{Matrix Form}

We can write the system as matrices $A, B, C$:

\[
A =
\begin{bmatrix}
1 & 2 & -1 & 0 & 0 & 0 & 0 & 0 \\
0 & 1 & 3 & 0 & 0 & 0 & 0 & 0 \\
0 & 0 & 0 & 1 & 0 & 0 & 0 & 0 \\
0 & 0 & 0 & 0 & 1 & 0 & 0 & 0 \\
3 & 0 & 0 & 0 & 0 & 4 & -2 & 0
\end{bmatrix},
\]

\[
B =
\begin{bmatrix}
1 & 0 & 0 & 0 & 0 & 0 & 0 & 0 \\
1 & 0 & 0 & 0 & 0 & 0 & 0 & 0 \\
0 & 0 & 0 & 1 & 0 & 0 & 0 & 0 \\
0 & 0 & 0 & 0 & 1 & 0 & 0 & 0 \\
1 & 0 & 0 & 0 & 0 & 0 & 0 & 0
\end{bmatrix},
\]

\[
C =
\begin{bmatrix}
0 & 0 & 0 & 1 & 0 & 0 & 0 & 0 \\
0 & 0 & 0 & 0 & 1 & 0 & 0 & 0 \\
0 & 0 & 0 & 0 & 0 & 1 & 0 & 0 \\
0 & 0 & 0 & 0 & 0 & 0 & 1 & 0 \\
0 & 0 & 0 & 0 & 0 & 0 & 0 & 1
\end{bmatrix}.
\]

\subsection{Verification Example}

Let
\[
x_1 = 2, \quad x_2 = 1.
\]

Then
\[
z_1 = 4, \quad z_2 = 5,
\]
\[
h_1 = 16, \quad h_2 = 25,
\]
\[
y = 17.
\]

So the witness is
\[
w = (1, 2, 1, 4, 5, 16, 25, 17).
\]

Check one constraint:
\[
\langle A_5, w \rangle = 3 + 4(16) - 2(25) = 17,
\quad
\langle B_5, w \rangle = 1,
\quad
\langle C_5, w \rangle = 17.
\]

Thus
\[
17 \cdot 1 = 17.
\]

\subsection{General Construction Recipe}

Given any feed-forward model:

\begin{enumerate}
\item Quantize inputs and parameters.
\item Introduce a variable for each intermediate value.
\item Decompose computation into primitive operations.
\item Encode affine operations as
\[
(\text{linear form}) \cdot 1 = \text{output}.
\]
\item Encode multiplications as
\[
a \cdot b = c.
\]
\item Assemble all constraints into $(A, B, C)$.
\item Provide a witness containing all values.
\end{enumerate}

\subsection{Remarks}

\begin{itemize}
\item Neural networks operate over $\mathbb{R}$, while R1CS is defined over a finite field.
\item Therefore, zkML typically proves correctness of a \emph{quantized} version of the model.
\item Nonlinearities such as ReLU require special treatment (approximation, lookup, or custom constraints).
\end{itemize}

\medskip

\noindent
\textbf{Key Insight.}  
A neural network is a compositional function.  
R1CS expresses the same computation as a system of polynomial constraints.